
\usepackage{fullpage}
%\usepackage{enumitem}
\usepackage{paralist}

\usepackage{amsmath,amsthm,amsfonts,amssymb,mathtools,latexsym}
\usepackage[dvipsnames]{xcolor}
\usepackage{graphicx,hyperref,float}
\usepackage[ruled,vlined,linesnumbered]{algorithm2e}

\usepackage{tikz}
\usetikzlibrary{trees,shapes.geometric,positioning,arrows.meta}

\usepackage{listings}

  \definecolor{customgreen}{rgb}{0,0.6,0}
  \definecolor{customgray}{rgb}{0.5,0.5,0.5}
  \definecolor{custommauve}{rgb}{0.6,0,0.8}
  
  \lstset{ 
    basicstyle=\ttfamily\footnotesize,        % the size of the fonts that are used for the code
    breaklines=true,                 % sets automatic line breaking
    commentstyle=\color{customgreen},    % comment style
    firstnumber=1,                % start line enumeration with line 1000
    frame=single,	                   % adds a frame around the code
    keepspaces=true,                 % keeps spaces in text, useful for keeping indentation of code (possibly needs columns=flexible)
    keywordstyle=\color{blue},       % keyword style
    numbers=left,                    % where to put the line-numbers; possible values are (none, left, right)
    numbersep=10pt,                   % how far the line-numbers are from the code
    numberstyle=\tiny\color{customgray}, % the style that is used for the line-numbers
    rulecolor=\color{black},         % if not set, the frame-color may be changed on line-breaks within not-black text (e.g. comments (green here))
    showspaces=false,                % show spaces everywhere adding particular underscores; it overrides 'showstringspaces'
    showstringspaces=false,          % underline spaces within strings only
    showtabs=false,                  % show tabs within strings adding particular underscores
    stepnumber=1,                    % the step between two line-numbers. If it's 1, each line will be numbered
    stringstyle=\color{custommauve},     % string literal style
    tabsize=2,	                    % sets default tabsize to 2 spaces
    title=\lstname                   % show the filename of files included with \lstinputlisting; also try caption instead of title
  }


\usepackage{array}
\usepackage{multicol}
\usepackage{booktabs}

\usepackage{hyperref}
\usepackage{cleveref}

\usepackage[ruled,vlined,linesnumbered]{algorithm2e}






\definecolor{oxblue}{RGB}{0, 30, 150}



% box around soln
\newcommand{\answerbox}[1]{\fbox{\textcolor{white}{\rule{#1}{24pt}}}}
\newcommand{\answerboxshort}[1]{\fbox{\textcolor{white}{\rule{#1}{14pt}}}}
\newcommand{\answerboxtall}[2]{\fbox{\textcolor{white}{\rule{#1}{#2}}}}


\newcommand\mycommfont[1]{\footnotesize\ttfamily\textcolor{blue}{#1}}
\SetCommentSty{mycommfont}



\pdfpagewidth 8.5in
\pdfpageheight 11in 

\oddsidemargin 0in
\evensidemargin 0in


\newtheorem{claim}{Claim}
\newtheorem{definition}{Definition}
\newtheorem{theorem}{Theorem}
\newtheorem{lemma}{Lemma}
\newtheorem{observation}{Observation}
\newtheorem{question}{Question}
\newtheorem{problem}{Problem}

\newcommand{\re}{{\mathbb{R}}}
\newcommand{\floor}[1]{\lfloor {#1} \rfloor}
\newcommand{\ceil}[1]{\lceil {#1} \rceil}
\newcommand{\paren}[1]{\left( {#1} \right)}
\newcommand{\bparen}[1]{\big( {#1} \big)}
\newcommand{\Bparen}[1]{\Big( {#1} \Big)}

\newcommand{\set}[1]{\left\{ {#1} \right\}}

\DeclareMathOperator{\E}{\mathbb{E}}
\DeclareMathOperator{\StdDev}{\mathbb{\text{\rm StdDev}}}


\newcommand{\adamnote}[1]{\marginpar{{\tiny \sf \color{orange} Adam: {#1}}}}


\newenvironment{solution}
  {\medskip\noindent\color{oxblue}}
  {\medskip}

\ifnum\showsolutions=0
    \newsavebox{\discardbox}
    \renewenvironment{solution}{\setbox\discardbox=\vbox\bgroup}
    {\egroup}
\fi 

